\section{Threats to Validity}
\label{sec:threats_to_validity}

Following the guidelines of Wohlin et al. \cite{wohlin_2012}, we identify potential threats to validity and the mitigation strategies implemented in our experimental design.

\subsection{Construct Validity}
Construct validity concerns the degree to which our experimental metrics accurately reflect the theoretical constructs under study. To prevent measurement bias, the primary dependent variables ($Acc$ and $MTTDiag$) were collected through an automated, blind execution harness that recorded authoritative UTC timestamps directly from the Kubernetes API Server and telemetry query engines. Root-cause isolation verdicts were mapped strictly against the objective ground truth ($\text{target}_i$) injected by the Chaos Mesh controller, eliminating human observer subjectivity. Furthermore, our $300\,\text{s}$ timeout penalty for misdiagnoses and unresolved incidents prevents right-censored data from artificially flattering diagnostic latency metrics.

\subsection{Internal Validity}
Internal validity addresses whether observed diagnostic improvements and overhead measurements were caused by the PANOPTES architecture rather than confounding environmental factors. To eliminate confounding memory thrashing and CPU starvation, all worker nodes were configured with hard Linux cgroup resource limits ($2.0\,\text{vCPUs}$, $4096\,\text{MiB}$ RAM). Following our Markdown Architectural Decision Record \textbf{ADR-0005}, non-transactional auxiliary workloads (e.g., demo AI agents and search indexers) were scaled to zero replicas, preserving cluster stability and guaranteeing identical baseline conditions across all campaigns.

To mitigate pseudoreplication \cite{hurlbert_1984}, steady-state overhead sampling was conducted at $\Delta t = 16\,\text{s}$ intervals, intentionally exceeding the $15\,\text{s}$ Kubelet scraper and metric aggregation window to guarantee independent observations. Finally, cluster state, Prometheus metrics, and temporary pods were purged and verified via pre-flight health scripts between successive fault campaigns.

\subsection{External Validity}
External validity concerns the generalizability of our findings to real-world production environments. While our evaluation was conducted on an isolated K3s/K3d testbed rather than a multi-region managed cloud cluster (e.g., EKS, GKE), the CNCF OpenTelemetry Astronomy Shop benchmark was explicitly selected because its polyglot microservice architecture, mixed synchronous/asynchronous IPC (HTTP, gRPC, Apache Kafka), and realistic Locust user traffic faithfully reproduce production cloud-native complexities. Moreover, because PANOPTES relies strictly on vendor-neutral CNCF standards (OpenTelemetry, Prometheus, W3C TraceContext), the architectural principles and algorithms are directly portable to any standard Kubernetes distribution.

\subsection{Conclusion Validity}
Conclusion validity relates to the statistical rigor of the inferences drawn from the experimental data. Given that latency and overhead distributions in distributed systems typically exhibit severe skewness and multimodal tails, we avoided assuming normal distributions and relied exclusively on non-parametric statistical tests: Fisher's Exact Test and McNemar's Paired Test for categorical diagnostic accuracy, and the Wilcoxon Signed-Rank Test along with the Vargha-Delaney $\hat{A}_{12}$ effect size for diagnostic latency distributions. All primary hypothesis tests yielded significance levels well below the standard threshold ($\alpha = 0.05$), confirming that the empirical findings are robust.
